\chapter{How \app\ is verified}\label{app:verification}

\app\ computes engineering results, so it is tested against independent references rather
than against its own earlier output. An automated test suite, a few hundred tests, runs
in seconds and covers the engine, the file format, the command line and the window, which
is driven without a screen.

\section{References}

\begin{description}
  \item[MATLAB reference solutions.] Every numeric variable of the MATLAB solutions of
    the four assignments, and of the LF variant, is stored at full precision (17
    significant digits) for $A = 80$ and $A = 0$. The tests run every case of the example
    projects at both values of \code{A} and compare with a relative tolerance of
    $10^{-9}$, and $10^{-8}$ for the load flow. A refactoring that moves a result in its
    eighth digit fails.
  \item[The assignment statements.] The $A = 0$ check values printed in the statements are
    an independent check, not a comparison with anything \app\ produced, for example
    $Z_I = 22.5156 + j11.0767$~pu in S1 and $P_\text{loss} = 26.4240$~MW in LF.
  \item[pandapower.] The load flow is cross-checked against pandapower, an independent
    open-source solver, on the LF network: voltages, angles, slack power and losses agree
    to $10^{-8}$. This test is opt-in (\cli{PSA\_ORACLES=1}), because importing pandapower
    takes longer than the rest of the suite.
  \item[Finite differences.] The load-flow Jacobian is compared with a finite-difference
    Jacobian at a random state, and its diagonal with the textbook formulas of
    \cref{sec:nr}.
\end{description}

\section{Self-checks in every result}

The results carry their own cross-checks (\cref{ch:studies}): the two-port matrix by
cascade and by Kron reduction, the equivalent impedance by Kron reduction and through
the two-port, the Thevenin impedance by the impedance matrix and by Kron reduction, the
Thevenin voltage by the nodal solve and by the MATLAB route, $\det F = 1$ for reciprocal
networks, the symmetry of the admittance matrix, and the load-flow power balance. A
residual well above rounding level is a sign that something in the network is not what
you think.

\section{Properties that must hold}

Besides fixed numbers, the tests check relations any correct solver satisfies. A few
examples:
\begin{itemize}
  \item an unbalanced load with phase factors $(1, 1, 1)$ gives the balanced solution;
  \item the general fault solver reproduces the closed-form currents of the
    three-phase, line-to-ground, line-to-line and double line-to-ground faults;
  \item a \code{Dyn} transformer is a shunt on its star side in the zero sequence, and
    nothing on its delta side;
  \item two generators sharing a bus give the same flow as one generator with their sum;
  \item a constant-impedance load draws $S_n U^2$ at the solved voltage.
\end{itemize}

\section{Running the tests yourself}

From the repository's root:

\begin{terminal}
python -m pytest                            # the whole suite
python -m pytest tests/test_lf.py           # one file
PSA_ORACLES=1 python -m pytest tests/test_lf.py -k pandapower
python -m unittest discover -s tests -t .   # without pytest
\end{terminal}

The reference tables of \cref{ch:components,ch:studies} are generated from the program
itself, by \cli{tools/make\_manual.py}. A test fails when they fall out of step with the
code.
